\documentclass[11pt,a4paper]{article}

\usepackage[a4paper,margin=1in]{geometry}
\usepackage{amsmath,amssymb}
\usepackage{booktabs}
\usepackage{graphicx}
\usepackage{hyperref}
\usepackage{float}

\title{A Technical Note on an Apparently Unlisted Integer Sequence from Asymmetric Prime--Composite Neighborhoods}
\author{Aradhya Haldikar}
\date{\today}

\begin{document}

\maketitle

\begin{abstract}
We document an integer sequence defined by a simple asymmetric condition involving prime and composite numbers. For each integer \(n \ge 3\), let \(a(n)\) be the least positive integer \(k\) such that \(n-k\) is prime and \(n+k\) is composite. We provide initial terms, computational results up to \(n=100000\), summary statistics, graphical observations, and open questions. Searches of the On-Line Encyclopedia of Integer Sequences (OEIS) using initial terms returned no exact match at the time of writing.
\end{abstract}

\section{Introduction}

Many integer sequences arise from local prime behavior, such as prime gaps, nearest-prime distances, or counting functions involving primes in intervals. This note considers a simple asymmetric variant: the smallest symmetric offset from \(n\) for which the left side is prime and the right side is composite.

Although elementary to define, the resulting sequence displays frequent small values together with occasional larger spikes.

\section{Definition}

Let \(\mathbb{P}\) denote the set of prime numbers. For each integer \(n \ge 3\), define

\[
a(n)=\min \{k \ge 1 : n-k \in \mathbb{P},\; n+k \notin \mathbb{P},\; n+k>1\}.
\]

Equivalently, \(a(n)\) is the least positive integer \(k\) such that:

\begin{itemize}
    \item \(n-k\) is prime,
    \item \(n+k\) is composite.
\end{itemize}

\section{Initial Terms}

For \(n=3,4,5,\dots\), the sequence begins:

\[
1,2,3,3,2,1,6,5,4,9,2,1,10,5,4,7,2,1,4,3,4,1,2,7,8,5,6,19,
\]

\[
2,1,2,5,4,13,8,1,10,9,4,13,2,1,4,3,4,1,2,7,4,5,10,1,2,9,20,\dots
\]

\section{Examples}

\subsection*{Example 1: \(n=14\)}

\[
14-1=13 \text{ is prime}, \qquad 14+1=15 \text{ is composite}.
\]

Hence:

\[
a(14)=1.
\]

\subsection*{Example 2: \(n=30\)}

At \(k=19\),

\[
30-19=11 \text{ is prime}, \qquad 30+19=49 \text{ is composite}.
\]

No smaller positive \(k\) satisfies both conditions, so:

\[
a(30)=19.
\]

\section{Computational Results}

Direct computation verified that \(a(n)\) exists for all tested integers

\[
3 \le n \le 100000.
\]

No exceptions were found in this range.

\subsection*{Summary Statistics}

\begin{center}
\begin{tabular}{lc}
\toprule
Quantity & Value \\
\midrule
Mean of \(a(n)\) & 11.0621 \\
Median of \(a(n)\) & 8 \\
Minimum of \(a(n)\) & 1 \\
Maximum of \(a(n)\) & 134 \\
Standard deviation & 10.2241 \\
\bottomrule
\end{tabular}
\end{center}

The empirical mean exceeds the median, indicating a right-skewed distribution with many small values and occasional larger spikes.

\section{Graphs}

\begin{figure}[H]
\centering
\includegraphics[width=0.85\textwidth]{Scatterplot.png}
\caption{Scatter plot of \(a(n)\) for \(3 \le n \le 100000\).}
\end{figure}

\begin{figure}[H]
\centering
\includegraphics[width=0.85\textwidth]{Histogram.png}
\caption{Histogram of computed values of \(a(n)\) up to \(100000\).}
\end{figure}

\section{Observed Behavior}

Empirically:

\begin{enumerate}
    \item Small values occur frequently.
    \item Many even integers satisfy \(a(n)=1\), especially when \(n-1\) is prime and \(n+1\) is composite.
    \item Larger values occur irregularly.
    \item The maximum observed value up to \(100000\) is \(134\).
\end{enumerate}

\section{OEIS Search Note}

An OEIS search using several initial blocks of terms returned no exact match at the time of writing. This does not exclude equivalent formulations already present under a different description.

\section{Algorithm}

A direct algorithm is:

\begin{verbatim}
for n = 3,4,5,...
    for k = 1,2,3,...
        if isPrime(n-k) and isComposite(n+k):
            output k
            break
\end{verbatim}

\section{Code Availability}

Source code, computed data, and manuscript files are available at:

\begin{center}
\url{https://github.com/Dhpla12/prime-composite-asymmetric-sequence}
\end{center}

\section{Open Questions}

\begin{enumerate}
    \item Does \(a(n)\) exist for every integer \(n \ge 3\)?
    \item Is the sequence unbounded?
    \item What is the asymptotic average order of \(a(n)\)?
    \item What proportion of integers satisfy \(a(n)=1\)?
    \item How are record values related to prime gaps?
\end{enumerate}

\section{Conclusion}

We documented an apparently unlisted sequence arising from asymmetric prime--composite neighborhoods around integers. The sequence is easy to compute, empirically well-defined for a large tested range, and exhibits nontrivial local variation.

\begin{thebibliography}{9}

\bibitem{hardywright}
G. H. Hardy and E. M. Wright,
\textit{An Introduction to the Theory of Numbers},
Oxford University Press.

\bibitem{oeis}
The On-Line Encyclopedia of Integer Sequences,
\url{https://oeis.org}

\end{thebibliography}

\end{document}
